En un experiment per a mesurar la força magnètica sobre una intensitat de corrent, situem un imant sobre una balança. L'imant té una llargària $L = 20,0\ \mathrm{cm}$. Entre els pols de l'imant hi passa un fil conductor de 15,0~cm de llargària que està orientat perpendicularment al camp magnètic. El camp magnètic creat per l'imant és homogeni en la regió on es troba el fil conductor. La balança marca 239~g quan no passa cap intensitat de corrent pel fil, mentre que quan hi circula una intensitat de 8,00~A la balança marca 243~g.

\figura[0.7]{fig1.png}

\begin{apartat}{1,25}
Quina és la intensitat del camp magnètic creat per l'imant? Justifiqueu per què la lectura de la balança augmenta quan hi passa corrent.
\end{apartat}

\begin{apartat}{1,25}
Posteriorment, traiem l'imant i fem passar la mateixa intensitat de corrent, 8,00~A, per un fil recolzat sobre la balança, paral·lel al fil anterior i situat a una distància $r$ per sota.

\figura[0.22]{fig2.png}

Quina hauria de ser aquesta distància $r$ per a obtenir la mateixa força sobre el primer fil conductor?
\end{apartat}

\blocetiquetat{Dada:}{La intensitat del camp magnètic creat per un conductor rectilini pel qual circula un corrent $I$ en un punt situat a una distància $r$ del conductor és: $B = \dfrac{\mu_0\, I}{2\pi\, r}$, en què $\mu_0 = 4\pi\times 10^{-7}\ \mathrm{T\,m\,A^{-1}}$.}
